\chapter{定向前置}
\label{ch:prerequisites}

\section{怎样使用本章}
\label{sec:ch00-how-to-use}

本章只放后文实际调用的最低充分前置。请把内容分为“第一章前必读”“到指定章节前即时补读”和“仅在证明中查阅”，不要要求读者一次学完所有工具。

\section{前置概念模板}
\label{sec:ch00-concept}

\begin{definition}
在此给出完整定义。
\end{definition}

\begin{note}
解释该定义解决什么问题、后文在哪里调用，以及第一遍需要掌握到什么程度。
\end{note}

% 可选前置自检片段：只在确有必要时取消注释，按实际前置与阅读时机改题；
% 这里只放无答案题面。对应答案片段在 appendices/solutions.tex。
% \section{前置自检}
% \label{sec:ch00-selfcheck}
% \noindent\textbf{第一章前：}回答下题；若改为其他阅读时机，请同时修改此处。
% \begin{enumerate}
%   \item\label{check:ch00-01}
%   设线性映射 $T\colon V\to V$。写出 $T$-不变子空间的定义，
%   并说明 $\ker T$ 是否 $T$-不变。\quad
%   \hyperref[checksol:ch00-01]{查看参考答案}
% \end{enumerate}

